\documentclass[11pt]{article}

\usepackage[margin=1in]{geometry}
\usepackage{listings}
\usepackage{xcolor}
\usepackage{longtable}
\usepackage{booktabs}
\usepackage{array}
\usepackage{hyperref}
\usepackage{tabularx}
\usepackage{caption}

\hypersetup{
    colorlinks=true,
    linkcolor=blue,
    urlcolor=blue,
}

\definecolor{codegray}{rgb}{0.95,0.95,0.95}
\definecolor{codecomment}{rgb}{0.25,0.5,0.25}
\definecolor{codekeyword}{rgb}{0.0,0.0,0.7}
\definecolor{codestring}{rgb}{0.6,0.1,0.1}

\lstdefinestyle{cppstyle}{
    backgroundcolor=\color{codegray},
    commentstyle=\color{codecomment}\itshape,
    keywordstyle=\color{codekeyword}\bfseries,
    stringstyle=\color{codestring},
    basicstyle=\ttfamily\footnotesize,
    breakatwhitespace=false,
    breaklines=true,
    captionpos=b,
    keepspaces=true,
    numbers=left,
    numbersep=5pt,
    showspaces=false,
    showstringspaces=false,
    showtabs=false,
    tabsize=2,
    frame=single,
    language=C++,
}

\lstdefinestyle{pystyle}{
    backgroundcolor=\color{codegray},
    commentstyle=\color{codecomment}\itshape,
    keywordstyle=\color{codekeyword}\bfseries,
    stringstyle=\color{codestring},
    basicstyle=\ttfamily\footnotesize,
    breakatwhitespace=false,
    breaklines=true,
    captionpos=b,
    keepspaces=true,
    numbers=left,
    numbersep=5pt,
    showspaces=false,
    showstringspaces=false,
    showtabs=false,
    tabsize=2,
    frame=single,
    language=Python,
}

\lstdefinestyle{inistyle}{
    backgroundcolor=\color{codegray},
    basicstyle=\ttfamily\footnotesize,
    breaklines=true,
    captionpos=b,
    keepspaces=true,
    numbers=left,
    numbersep=5pt,
    frame=single,
}

\title{EE120 Final Project Report \\
       \large Networked RGB Lamp System (Arduino + ESP8266)}
\author{Gabriel Jordaan, Sam Schultz, Nick Sankovich}
\date{May 2026}

\begin{document}

\maketitle

\section{Project Overview}

This project implements a network of four RGB lamps that communicate via a
mesh of ESP8266 Wi-Fi nodes paired with Arduino Uno boards. Each lamp has a
sensor that, when activated, signals its adjacent lamps over the network to
turn on for a fixed duration. Each lamp consists of two microcontrollers:
an Arduino Uno (sensor / RGB driver / coordinator) and an ESP8266 (Wi-Fi
transport, with one node configured as a router/access-point and the rest as
clients).

\section{Source Code}

\subsection{Arduino: \texttt{ardu code/src/main.cpp}}
\lstinputlisting[style=cppstyle]{ardu code/src/main.cpp}

\subsection{Arduino: \texttt{lib/Internet/src/Internet.h}}
\lstinputlisting[style=cppstyle]{ardu code/lib/Internet/src/Internet.h}

\subsection{Arduino: \texttt{lib/Internet/src/Internet.cpp}}
\lstinputlisting[style=cppstyle]{ardu code/lib/Internet/src/Internet.cpp}

\subsection{Arduino: \texttt{lib/SignalCoordinator/src/SignalCoordinator.h}}
\lstinputlisting[style=cppstyle]{ardu code/lib/SignalCoordinator/src/SignalCoordinator.h}

\subsection{Arduino: \texttt{lib/SignalCoordinator/src/SignalCoordinator.cpp}}
\lstinputlisting[style=cppstyle]{ardu code/lib/SignalCoordinator/src/SignalCoordinator.cpp}

\subsection{Arduino: \texttt{lib/Light/src/Light.h}}
\lstinputlisting[style=cppstyle]{ardu code/lib/Light/src/Light.h}

\subsection{Arduino: \texttt{lib/Light/src/Light.cpp}}
\lstinputlisting[style=cppstyle]{ardu code/lib/Light/src/Light.cpp}

\subsection{Arduino: \texttt{lib/Sensor/src/Sensor.h}}
\lstinputlisting[style=cppstyle]{ardu code/lib/Sensor/src/Sensor.h}

\subsection{Arduino: \texttt{lib/Sensor/src/Sensor.cpp}}
\lstinputlisting[style=cppstyle]{ardu code/lib/Sensor/src/Sensor.cpp}

\subsection{Arduino: \texttt{lib/Utils/src/Utils.h}}
\lstinputlisting[style=cppstyle]{ardu code/lib/Utils/src/Utils.h}

\subsection{Arduino: \texttt{lib/Utils/src/Utils.cpp}}
\lstinputlisting[style=cppstyle]{ardu code/lib/Utils/src/Utils.cpp}

\subsection{Arduino: \texttt{platformio.ini}}
\lstinputlisting[style=inistyle]{ardu code/platformio.ini}

\subsection{ESP8266: \texttt{ESP\_code\_EE120\_final/src/main.cpp}}
\lstinputlisting[style=cppstyle]{ESP_code_EE120_final/src/main.cpp}

\subsection{ESP8266: \texttt{platformio.ini}}
\lstinputlisting[style=inistyle]{ESP_code_EE120_final/platformio.ini}

\subsection{Build / Macro Patcher: \texttt{macrocompile.py}}
\lstinputlisting[style=pystyle]{macrocompile.py}

\section{Pin Input / Output Table}

The table below lists every pin used across the Arduino Uno and the
ESP8266 NodeMCU, the component (software module) that owns the pin, and
the role that pin plays in the system.

\begin{longtable}{>{\raggedright\arraybackslash}p{0.06\textwidth}
                  >{\raggedright\arraybackslash}p{0.10\textwidth}
                  >{\raggedright\arraybackslash}p{0.07\textwidth}
                  >{\raggedright\arraybackslash}p{0.18\textwidth}
                  >{\raggedright\arraybackslash}p{0.45\textwidth}}
\toprule
\textbf{Board} & \textbf{Direction} & \textbf{Pin} & \textbf{Component} & \textbf{Function} \\
\midrule
\endfirsthead

\toprule
\textbf{Board} & \textbf{Direction} & \textbf{Pin} & \textbf{Component} & \textbf{Function} \\
\midrule
\endhead

Arduino & OUTPUT & D10 & Light (Red LED) & PWM drive for the red channel of the RGB lamp (\texttt{analogWrite}). \\
Arduino & OUTPUT & D11 & Light (Green LED) & PWM drive for the green channel of the RGB lamp. \\
Arduino & OUTPUT & D3  & Light (Blue LED) & PWM drive for the blue channel of the RGB lamp. \\
Arduino & INPUT  & A0  & Sensor (Analog) & Reads the analog motion / proximity sensor; lamp is considered active when the value crosses \texttt{SENSOR\_THRESHOLD} (512). \\
Arduino & INPUT  & D13 & Sensor (Digital) & Alternative digital sensor input (used when \texttt{SENSOR\_MODE = "DIGITAL"}). \\
Arduino & OUTPUT & D12 & Internet (TX-ADD) & Pulses HIGH for $\sim$1\,s to tell the local ESP8266 to broadcast an ``ADD'' request, marking adjacent lamps as active on the shared network list. \\
Arduino & OUTPUT & D5  & Internet (TX-REMOVE) & Pulses HIGH for $\sim$1\,s to tell the local ESP8266 to send a ``REMOVE'' acknowledgement, removing this lamp from the network active list. \\
Arduino & INPUT  & D6  & Internet (RX) & Reads the activation line from the local ESP8266; HIGH means this lamp's index appears in the network active list and the lamp should turn on. \\
\midrule
ESP8266 & INPUT  & GPIO4 (D2)  & Internet (RX-ADD) & Receives the Arduino's ``add adjacent lamps'' pulse; on a rising edge the ESP sends \texttt{ADD <n>} for each adjacent index over Wi-Fi. \\
ESP8266 & INPUT  & GPIO16 (D0) & Internet (RX-REMOVE) & Receives the Arduino's ``remove me'' pulse; the ESP sends \texttt{REMOVE <LIGHT\_INDEX>} over Wi-Fi to clear itself from the active list. \\
ESP8266 & OUTPUT & GPIO14 (D5) & Internet (TX) & Drives the Arduino's RX line HIGH whenever this lamp's \texttt{LIGHT\_INDEX} is present in the broadcast active list (i.e.\ this lamp should turn on). \\
\bottomrule
\caption{Pin assignments across the Arduino Uno and ESP8266 NodeMCU for one lamp node.}
\label{tab:pinout}
\end{longtable}

\section{Git Commits Tied to Each Feature}

The table below maps each major feature of the project to the git commits
that introduced or finalized it. Commit hashes are short SHAs from the
project repository.

\begin{longtable}{>{\raggedright\arraybackslash}p{0.28\textwidth}
                  >{\raggedright\arraybackslash}p{0.20\textwidth}
                  >{\raggedright\arraybackslash}p{0.45\textwidth}}
\toprule
\textbf{Feature} & \textbf{Commit(s)} & \textbf{Description} \\
\midrule
\endfirsthead

\toprule
\textbf{Feature} & \textbf{Commit(s)} & \textbf{Description} \\
\midrule
\endhead

Project skeleton / PlatformIO scaffold &
\texttt{8fd7f04} \newline (Initial commit) &
Created the PlatformIO project layout, \texttt{.gitignore}, and stub \texttt{src/main.cpp}. \\

Module skeletons (Internet, Light, Sensor, SignalCoordinator) &
\texttt{c5e277e} \newline (skeleton) &
Stubbed out the \texttt{Internet}, \texttt{Light}, \texttt{Sensor} and \texttt{SignalCoordinator} libraries with empty class methods. \\

Skeleton wiring \& \texttt{Utils} helper &
\texttt{92fed2d} \newline (Skeleton done) &
Connected the modules into \texttt{main.cpp} via \texttt{begin()}/\texttt{update()}, added the \texttt{Utils} library, and fleshed out an early \texttt{SignalCoordinator}. \\

RGB pin macros &
\texttt{8542f63} \newline (Put RGB macros in Light.h) &
Centralized the \texttt{RED\_PIN}, \texttt{GREEN\_PIN}, \texttt{BLUE\_PIN} pin numbers in \texttt{Light.h}. \\

RGB lamp output (Light module) &
\texttt{144188f} \newline (Completed Light.cpp) &
Implemented \texttt{Light::begin} (\texttt{pinMode} OUTPUT) and \texttt{Light::setColor} using \texttt{analogWrite} on the three PWM pins. \\

Repository restructure into Arduino + ESP subprojects &
\texttt{2c1c8f3} \newline (Added ESP subrepo) &
Split the codebase into \texttt{ardu code/} and \texttt{ESP\_code\_EE120\_final/} so the Uno and ESP8266 firmware can be built as separate PlatformIO projects. \\

Cross-platform macro patching tool &
\texttt{109c159} \newline (Added gabage slop macro patching) &
Added \texttt{macrocompile.py}, which rewrites \texttt{\#define} values across both subprojects from a single source of truth so pin numbers stay in sync between the Arduino and ESP code. \\

ESP8266 Wi-Fi networking (router/client, ADD/REMOVE/LIST protocol) &
\texttt{38517a6} \newline (ESP CODE) &
Implemented the ESP firmware: SoftAP router mode with \texttt{WiFiServer}/client list, station-mode clients, and the \texttt{ADD <n>} / \texttt{REMOVE <n>} / \texttt{LIST ...} text protocol. Also added the matching Arduino-side \texttt{txMessage} / \texttt{rxMessage} / \texttt{txDeleteMessage} signal lines. \\

Analog sensor mode &
\texttt{eb52ebb}, \texttt{2b9086c} \newline (Update Sensor.cpp) &
Added analog sensor support with a configurable threshold and active-high / active-low logic; later tweaked the comparison direction. \\

Per-lamp adjacency macros &
\texttt{24aa73e} \newline (Pity macros for Nick S.) &
Added the \texttt{LAMP\_INDEX}-driven adjacency / router selection to \texttt{macrocompile.py} (each lamp's neighbours are auto-set), and matching ESP-side \texttt{ADJACENT}/\texttt{LIGHT\_INDEX} macros. \\

End-to-end integration polish &
\texttt{dd32614} \newline (Finishing touches) &
Final wiring of the Arduino main loop: TX-add / TX-remove pulse timing, the 10-second active timer, and the ack-back ``remove me'' message after receiving an active signal. \\

Final flashed configuration &
\texttt{6446a32} \newline (Flashed final 3) &
Last on-device tuning pass before demo --- pin number cleanup, sensor calibration, and merged \texttt{Internet} / \texttt{SignalCoordinator} edits used in the final flashed image. \\

\bottomrule
\caption{Mapping of features to the commits that introduced or finalized them.}
\label{tab:commits}
\end{longtable}

\end{document}
